\section{The quantitative pair-extension criterion}

The mate-line decomposition now turns the extension problem into one count:
empty fixed lines are the carriers, and nonfixed secant orbits are the only
first-order obstructions on them.

\begin{theorem}[Quadratic pair-extension criterion]
\label{thm:pair-criterion}
Let $\cC$ be a $\phi$-invariant $k$-arc in $\PG(2,s^2)$ with profile
$k=f+2e$.  With $E$, $N$, and $M$ as in
\eqref{eq:E} and \eqref{eq:M},
\begin{equation}\label{eq:uniform-lower}
  \Npair(\cC)\ge E\,\pospart{N-M}
  =E\,\pospart{\frac{s^2-s}{2}-fe-e(e-1)}.
\end{equation}
In particular, if $E>0$ and $M<N$, then $\cC$ admits a legal conjugate-pair
extension.
\end{theorem}

\begin{proof}
Fix $\ell\in\cE$.  It carries exactly $N$ candidate conjugate pairs.  A
fixed old secant meets $\ell$ in a fixed point and destroys no nonfixed
candidate.  Consider a nonfixed old secant orbit $\{m,\phi(m)\}$.  Unless
$m\cap\ell$ is fixed, the two intersections $m\cap\ell$ and
$\phi(m)\cap\ell$ form one conjugate candidate pair; thus the orbit destroys
at most one candidate.  There are $M$ such secant orbits, so at least
$\pospart{N-M}$ candidates survive on $\ell$.

A survivor lies on no old secant, and its empty mate line contains no point
of $\cC$.  Lemma~\ref{lem:pair-test} makes it legal.  Candidate sets on
distinct fixed lines are disjoint because a nonfixed pair has a unique mate
line.  Summing over the $E$ carriers proves \eqref{eq:uniform-lower}.
\end{proof}

\begin{corollary}[The Clebsch hexagon over the quadratic extension]
\label{cor:clebsch-pair-extension}
Let $H\subseteq\PG(2,11)$ be a Clebsch hexagon, viewed as a
Frobenius-invariant six-arc in $\PG(2,11^2)$
\cite{RuddClebschRigidity2026}.  Then
\[
  \Npair(H)=76\cdot55=4180.
\]
Consequently every invariant eight-arc $H\cup Q$ obtained by adjoining one
such pair has exactly $4179$ alternate conjugate-pair repairs after $Q$ is
deleted.  In coding language, the scalar extension of the Clebsch
$[6,3,4]_{11}$ code has $4180$ Frobenius-paired two-column extensions to
$[8,3,6]_{121}$ MDS codes.
\end{corollary}

\begin{proof}
Here $(s,f,e)=(11,6,0)$.  There are
\[
 E=11^2+11+1-\left(6\cdot12-\binom62\right)=76
\]
fixed lines disjoint from $H$, and each contains
$N=(11^2-11)/2=55$ nonfixed conjugate pairs.  Every old secant is fixed, so
$M=0$ and every one of these pairs is legal.  Conversely, the mate line of
any legal nonfixed pair is a fixed line disjoint from $H$, which proves
equality rather than only the lower bound.  Excluding the selected pair
$Q$ leaves $4180-1$ alternatives.
\end{proof}

\begin{remark}
The count in Corollary~\ref{cor:clebsch-pair-extension} uses only that $H$
is a six-arc rational over $\F_{11}$; it is not another characterization of
the Clebsch class.  Its role is to show how the exceptional code studied in
\cite{RuddClebschRigidity2026} enters the equivariant extension and repair
theory developed here.
\end{remark}

The preceding proof deliberately uses only the uniform upper bound $M$ on
the number of forbidden candidates per carrier.  Retaining the actual
forbidden support gives both a sharper bound and an exact identity.

For $\ell\in\cE$, let $\cQ_\ell$ be its $N$ candidate pairs and let
$\cF_\ell\subseteq\cQ_\ell$ be the pairs hit by at least one nonfixed old
secant orbit.  Then
\begin{equation}\label{eq:heterogeneous}
  \Npair(\cC)=\sum_{\ell\in\cE}
  \left(|\cQ_\ell|-|\cF_\ell|\right).
\end{equation}
The same formula applies when carrier capacities vary.  In particular, if
$|\cF_\ell|\le U_\ell$, then
\[
  \Npair(\cC)\ge
  \sum_{\ell\in\cE}\pospart{|\cQ_\ell|-U_\ell}.
\]
